% Manual técnico de la solución base — Proyecto Oráculo (ADAG1), Grupo G01
% Compilar con pdfLaTeX (por ejemplo en Overleaf). Dos pasadas para el índice.
\documentclass[12pt,letterpaper]{article}

\usepackage[utf8]{inputenc}
\usepackage[T1]{fontenc}
\usepackage[spanish,es-noshorthands,es-tabla,es-nodecimaldot]{babel}
\usepackage{newtxtext,newtxmath}
\usepackage[margin=2.54cm]{geometry}
\usepackage{setspace}
\usepackage{amsmath,amsthm}
\usepackage{booktabs,tabularx,array}
\usepackage{float}
\usepackage[labelfont=it,labelsep=period,justification=raggedright,singlelinecheck=false]{caption}
\usepackage{enumitem}
\usepackage{titlesec}
\usepackage{algorithm}
\usepackage[noend]{algpseudocode}
\usepackage{dirtree}
\usepackage{fancyhdr}
\usepackage[hidelinks]{hyperref}
\usepackage{xurl}

\onehalfspacing
\setlength{\parindent}{0pt}
\setlength{\parskip}{6pt}

% Todos los títulos en 12 pt, negrita
\titleformat{\section}{\normalfont\normalsize\bfseries}{\thesection.}{0.5em}{}
\titleformat{\subsection}{\normalfont\normalsize\bfseries}{\thesubsection.}{0.5em}{}
\titleformat{\subsubsection}{\normalfont\normalsize\bfseries\itshape}{}{0em}{}
\titlespacing*{\section}{0pt}{18pt}{6pt}
\titlespacing*{\subsection}{0pt}{12pt}{6pt}
\titlespacing*{\subsubsection}{0pt}{10pt}{4pt}
\setcounter{secnumdepth}{2}

\captionsetup[table]{position=top,skip=4pt}
\renewcommand{\arraystretch}{1.2}
\newcolumntype{L}{>{\raggedright\arraybackslash}X}

\pagestyle{fancy}
\fancyhf{}
\fancyfoot[R]{\thepage}
\renewcommand{\headrulewidth}{0pt}

% Pseudocódigo en español
\floatname{algorithm}{Algoritmo}
\algrenewcommand\algorithmicrequire{\textbf{Entrada:}}
\algrenewcommand\algorithmicensure{\textbf{Salida:}}
\algrenewcommand\algorithmicfor{\textbf{para}}
\algrenewcommand\algorithmicforall{\textbf{para cada}}
\algrenewcommand\algorithmicdo{\textbf{hacer}}
\algrenewcommand\algorithmicif{\textbf{si}}
\algrenewcommand\algorithmicthen{\textbf{entonces}}
\algrenewcommand\algorithmicreturn{\textbf{devolver}}
\algrenewcommand\alglinenumber[1]{\normalsize #1}

\theoremstyle{definition}
\newtheorem{proposicion}{Proposición}
\renewcommand{\proofname}{Demostración}

\newcommand{\Acc}{\operatorname{Acc}}

\begin{document}

% ---------------------------------------------------------------- portada
\begin{titlepage}
\centering
\vspace*{1cm}
{\bfseries Universidad de Caldas}\par
Ingeniería de Sistemas y Computación\par
Análisis y Diseño de Algoritmos\par
\vspace{4cm}
{\bfseries Manual técnico de la solución base}\par
Proyecto Oráculo (ADAG1), Grupo G01\par
Primera entrega\par
\vspace{4cm}
Autores:\par
Andrés Santiago Restrepo Duque\par
Keyla Daniela Cartagena Díaz\par
\vfill
Manizales\par
Noviembre 6 de 2026\par
\end{titlepage}

\renewcommand{\contentsname}{\centering Contenido}
\tableofcontents
\newpage

% ---------------------------------------------------------------- 1
\section{Introducción}

\subsection{Propósito del documento}
Este manual describe la construcción y el funcionamiento de la solución base del Proyecto Oráculo, de manera que una persona ajena al equipo pueda instalarla, ejecutarla, interpretar sus resultados y reproducir la configuración reportada en entrega.json sin explicaciones adicionales de los autores. Complementa el informe académico de la primera entrega, que contiene la formalización del problema y la demostración de NP-dureza; aquí se documenta la implementación.

\subsection{Alcance de la solución base}
La solución base resuelve el problema de seleccionar una configuración de prompt, es decir, un índice para cada una de las cinco ranuras del catálogo (rol, estrategia, formato, estilo y cierre) junto con una temperatura de generación, que maximice la precisión del oráculo sobre las 100 instancias de búsqueda. La precisión de una configuración es la fracción de instancias en las que la respuesta del modelo cumple todas sus restricciones verificables.

El método implementado es un ascenso por coordenadas con reinicios, al que se añade una regla de poda que utiliza las trazas de fallos del oráculo para descartar opciones. La solución entrega la mejor configuración observada durante la búsqueda; no se afirma que sea un óptimo global. Las estrategias de reparto de presupuesto del enunciado general, la exploración de temperaturas distintas de 0.0 y la validación experimental del crecimiento de $T(n)$ quedan fuera de esta versión y se listan en la sección~\ref{sec:pendiente}.

\subsection{Convenciones}
Los nombres de archivos, funciones y variables se escriben tal como aparecen en el código. Las rutas de la forma /content/\dots{} corresponden al sistema de archivos de Google Colab. Las cantidades en cursiva se definen en la sección~\ref{sec:notacion} y $|X|$ denota la cardinalidad del conjunto $X$.

% ---------------------------------------------------------------- 2
\section{Entorno de ejecución}

\subsection{Plataforma y hardware}
La generación de respuestas requiere una GPU compatible con CUDA. La solución se desarrolló y probó en Google Colab con una GPU NVIDIA T4 de 16~GB. La función cargar\_modelo de ayudas.py interrumpe la ejecución con un mensaje explícito si no detecta una GPU.

\begin{table}[H]
\caption{Requisitos de plataforma.}
\begin{tabularx}{\textwidth}{@{}p{4cm}L@{}}
\toprule
Elemento & Requisito \\
\midrule
Plataforma & Google Colab (entorno de ejecución con GPU) \\
GPU & NVIDIA T4, arquitectura sm\_75, 16 GB de VRAM \\
Sistema operativo & Linux de Colab (Ubuntu); no se probó en Windows ni macOS \\
Python & 3.10 o superior (versión provista por Colab) \\
Modelo de lenguaje & Qwen/Qwen3-1.7B (alias qwen17b), pesos en float16, aproximadamente 3,4 GB de VRAM \\
Almacenamiento & Google Drive, opcional, para conservar la caché entre sesiones \\
\bottomrule
\end{tabularx}
\end{table}

\subsection{Dependencias de software}
Las dependencias se declaran en requirements.txt y se instalan en la primera celda de cada notebook. La tabla~\ref{tab:dependencias} indica la función de cada una dentro de la solución.

\begin{table}[H]
\caption{Dependencias y su uso.}\label{tab:dependencias}
\begin{tabularx}{\textwidth}{@{}p{4.6cm}p{2.2cm}L@{}}
\toprule
Paquete & Versión & Uso en la solución \\
\midrule
torch & la de Colab & Inferencia del modelo en GPU \\
transformers & última & Carga del tokenizador y del modelo, generación de texto \\
accelerate & última & Ubicación del modelo en el dispositivo (device\_map) \\
bitsandbytes & $\geq$ 0.46.1 & Carga de modelos cuantizados a 4 bits (alias alternativos) \\
nltk, spacy, emoji, langdetect, immutabledict & última & Requeridos por los verificadores de restricciones de open-instruct \\
pandas, matplotlib & la de Colab & Tablas de evidencia y curva de búsqueda \\
tqdm & la de Colab & Barra de progreso del oráculo \\
\bottomrule
\end{tabularx}
\end{table}

\subsection{Recursos externos}
Además de los paquetes de Python, la ejecución descarga los siguientes recursos: el repositorio allenai/open-instruct, clonado en /content/open-instruct, que contiene el registro de verificadores IFEvalG usado por el oráculo; el modelo en\_core\_web\_sm de spaCy; los paquetes punkt, punkt\_tab, averaged\_perceptron\_tagger y averaged\_perceptron\_tagger\_eng de nltk, que se descargan manualmente porque el descargador de nltk falla detrás del proxy de Colab; y los pesos del modelo desde Hugging Face, sin autenticación.

\subsection{Consideraciones de la plataforma}
\begin{itemize}[leftmargin=*]
\item Después de la instalación es obligatorio reiniciar la sesión de Colab. Si no se reinicia, transformers puede quedar parcialmente actualizado y cargar\_modelo falla con un error poco descriptivo.
\item La T4 no ofrece bfloat16 nativo; ayudas.py fuerza float16 como tipo de cómputo.
\item Las sesiones gratuitas de Colab se desconectan por inactividad o por tiempo de uso. La caché del oráculo en Google Drive permite continuar una búsqueda interrumpida sin repetir generaciones (sección~\ref{sec:reanudacion}).
\item El repositorio debe ser público para que los notebooks puedan clonarlo desde Colab.
\end{itemize}

% ---------------------------------------------------------------- 3
\section{Arquitectura}

\subsection{Estructura del repositorio}
La figura~\ref{fig:arbol} muestra la organización de carpetas. Los archivos oraculo.py, ayudas.py, datos\_visibles.json y proyecto\_oraculo.ipynb provienen del repositorio del curso y no se modifican.

\begin{figure}[H]
\begin{singlespace}
\dirtree{%
.1 ADA-2026-2-GP01\_Andres\_y\_Keyla/.
.2 notebooks/.
.3 01\_busqueda.ipynb\DTcomment{búsqueda, trazas, validación y entrega.json}.
.3 02\_evaluacion.ipynb\DTcomment{evaluación de entrega.json}.
.2 resultados/\DTcomment{evidencia generada por 01\_busqueda.ipynb}.
.3 evaluaciones.csv, registro\_busqueda.json.
.3 derrotas.csv, poda.json, curva\_busqueda.png.
.3 ahorro\_poda.png, cache\_oraculo.json.gz.
.2 datos de\ \ prueba/\DTcomment{datos de prueba propios}.
.3 datos\_visibles\_100.json.
.2 pruebas/\DTcomment{scripts sin GPU}.
.3 prueba\_busqueda\_simulada.py\DTcomment{prueba con oráculo sintético}.
.3 comparar\_poda\_sintetica.py\DTcomment{búsqueda con y sin poda}.
.3 grafica\_ahorro\_poda.py\DTcomment{genera ahorro\_poda.png}.
.2 docs/.
.3 informe\_primera\_entrega.docx\DTcomment{informe académico}.
.3 manual\_tecnico.tex, manual\_tecnico.pdf\DTcomment{este manual}.
.2 entrega.json\DTcomment{configuración entregada}.
.2 oraculo.py, ayudas.py, datos\_visibles.json\DTcomment{archivos del curso}.
.2 proyecto\_oraculo.ipynb\DTcomment{notebook original del curso}.
.2 requirements.txt, README.md, .gitignore.
}
\end{singlespace}
\caption{Estructura de carpetas del repositorio.}\label{fig:arbol}
\end{figure}

\subsection{Responsabilidad de los módulos}

\begin{table}[H]
\caption{Módulos y responsabilidades.}
\begin{tabularx}{\textwidth}{@{}p{5cm}L@{}}
\toprule
Módulo & Responsabilidad \\
\midrule
oraculo.py & Arma el prompt de una configuración para cada instancia, genera la respuesta con el modelo, la verifica con open-instruct y guarda el resultado en caché. Se utiliza como caja negra. \\
ayudas.py & Carga el modelo y los datos, separa las particiones de búsqueda y validación, dibuja la curva de búsqueda y escribe entrega.json. \\
01\_busqueda.ipynb & Implementa la heurística, registra cada evaluación, aplica la poda basada en trazas, valida la configuración elegida y exporta la evidencia. \\
02\_evaluacion.ipynb & Comprueba el formato de entrega.json y mide su precisión sobre los datos de prueba o, en su ausencia, sobre la partición de validación. \\
prueba\_busqueda\_simulada.py & Ejecuta las funciones del notebook de búsqueda contra un oráculo sintético y verifica las propiedades de la sección~\ref{sec:propiedades}. \\
comparar\_poda\_sintetica.py & Ejecuta la misma búsqueda con la poda activada y desactivada sobre 50 configuraciones óptimas sintéticas y resume las consultas de cada variante. \\
grafica\_ahorro\_poda.py & Dibuja resultados/ahorro\_poda.png con las configuraciones evaluadas con y sin poda en los tres escenarios del informe. \\
datos\_visibles\_100.json & Datos de prueba propios del grupo, con el formato de datos\_visibles.json (sección~\ref{sec:datospropios}). \\
\bottomrule
\end{tabularx}
\end{table}

\subsection{Interfaces utilizadas}
La solución solo usa las interfaces públicas del código del curso, resumidas en la tabla~\ref{tab:interfaces}. Del catálogo de opciones se usa únicamente el número de opciones por ranura y el texto de una opción cuando se muestra la evidencia de una poda.

\begin{table}[H]
\caption{Interfaces públicas utilizadas.}\label{tab:interfaces}
\begin{tabularx}{\textwidth}{@{}p{5.6cm}LL@{}}
\toprule
Interfaz & Entrada & Salida \\
\midrule
Oraculo(modelo, datos, validacion, cache, lote, presupuesto) & Modelo cargado, particiones, ruta de caché y parámetros de lote & Objeto oráculo \\
oraculo.evaluar(config, instancias, semilla) & Configuración, lista de instancias, semilla & Resultado con precision, trazas y n \\
oraculo.validar(config, n, semilla) & Configuración y tamaño de muestra (None = todas) & Resultado sobre la partición de validación; imprime las familias más falladas \\
cargar\_modelo(alias) & Alias del modelo & Objeto Modelo (red, tokenizador, nombre) \\
cargar\_datos(ruta), dividir(datos) & Ruta del JSON; lista de instancias & Lista de instancias; tupla (busqueda, validacion) \\
entrega(grupo, config, semana, ruta) & Identificador, configuración, semana, ruta & Escribe entrega.json \\
\bottomrule
\end{tabularx}
\end{table}

\subsection{Flujo entre componentes}
El notebook de búsqueda llama a ayudas.py para obtener el modelo y las particiones, construye un objeto Oraculo y, en cada paso de la heurística, invoca oraculo.evaluar. El oráculo consulta primero su caché; solo genera respuestas para los pares de configuración e instancia que no estén guardados, y las verifica con el registro de open-instruct. La heurística recibe la precisión y las trazas, actualiza sus estructuras (sección~\ref{sec:estructuras}) y decide la siguiente configuración. Al terminar, el notebook llama a oraculo.validar y a entrega. El notebook de evaluación solo lee entrega.json y llama a oraculo.evaluar una vez.

% ---------------------------------------------------------------- 4
\section{Formatos de datos}

\subsection{Instancia}\label{sec:instancia}
Los archivos datos\_visibles.json y datos\_prueba.json contienen una lista JSON de instancias con los campos de la tabla~\ref{tab:instancia}.

\begin{table}[H]
\caption{Campos de una instancia.}\label{tab:instancia}
\begin{tabularx}{\textwidth}{@{}p{2.4cm}p{2.8cm}p{2.2cm}L@{}}
\toprule
Campo & Tipo & Obligatorio & Descripción \\
\midrule
id & entero & sí & Identificador único de la instancia \\
prompt & texto & sí & Petición completa al modelo, con sus restricciones redactadas \\
ids & lista de texto & sí & Familias de restricción de open-instruct, por ejemplo keywords:existence \\
kwargs & lista de objetos & sí & Parámetros de cada restricción, en el mismo orden que ids \\
restriccion & texto & no & Restricciones separadas por tabulador; la usa la opción de cierre que las repite \\
split & texto & no & busqueda, validacion o descartada \\
\bottomrule
\end{tabularx}
\end{table}

\subsection{Configuración}
Una configuración es un objeto con seis llaves. Las cinco ranuras toman un entero en $\{0, 1, \dots, 7\}$ y la temperatura toma un valor en $\{0.0, 0.3, 0.7\}$. Con la temperatura fija en 0.0 hay $8^5 = 32\,768$ configuraciones; con las tres temperaturas, $98\,304$. La solución base trabaja con temperatura 0.0. Por ejemplo:
\begin{center}
\{"rol": 0, "estrategia": 0, "formato": 0, "estilo": 0, "cierre": 0, "temperatura": 0.0\}
\end{center}

\subsection{Resultado del oráculo}
evaluar y validar devuelven un objeto con tres atributos: precision, número real en $[0, 1]$; n, número de instancias evaluadas; y trazas, lista que contiene solo las instancias falladas. Cada traza tiene id (instancia), violo (primera familia de restricción que no se cumplió) y salida (texto generado). Como la evaluación es de todo o nada, una instancia con varias restricciones incumplidas aparece una sola vez, con la primera que falló.

\subsection{Archivo entrega.json}
Contiene el identificador del grupo, la configuración elegida y la semana académica. Lo genera la función entrega al final del notebook de búsqueda y lo lee el notebook de evaluación:
\begin{center}
\{"grupo": "G01", "config": \{\dots\}, "semana": 5\}
\end{center}

\subsection{Archivos de evidencia}

\begin{table}[H]
\caption{Archivos generados en resultados/.}
\begin{tabularx}{\textwidth}{@{}p{4.4cm}L@{}}
\toprule
Archivo & Contenido \\
\midrule
evaluaciones.csv & Una fila por consulta nueva: número de consulta, reinicio, pasada, ranura modificada, opción, los cinco índices, precisión, aciertos y segundos \\
registro\_busqueda.json & Parámetros, resultado de cada reinicio, mejor configuración, precisión de búsqueda y de validación, fallos de validación por familia, opciones vivas al final y, por cada consulta, las familias violadas y los ids de las instancias falladas \\
derrotas.csv & Una fila por derrota observada: ranura, opción, vecindario, opción ganadora, precisiones, instancias perdidas y ganadas, reinicio y pasada \\
poda.json & Una entrada por opción descartada: ranura, opción, número de vecindarios, pérdida neta acumulada, evidencia de cada vecindario, familias violadas en las instancias perdidas y un extracto de la salida del modelo \\
curva\_busqueda.png & Precisión de cada consulta y mejor precisión acumulada \\
cache\_oraculo.json.gz & Copia comprimida de la caché del oráculo; permite rehacer el análisis sin generar de nuevo \\
\bottomrule
\end{tabularx}
\end{table}

% ---------------------------------------------------------------- 5
\section{Algoritmo de búsqueda}

\subsection{Notación}\label{sec:notacion}

\begin{table}[H]
\caption{Notación utilizada.}
\begin{tabularx}{\textwidth}{@{}p{4.6cm}L@{}}
\toprule
Símbolo & Significado \\
\midrule
$m$ & Número de ranuras; $m = 5$ \\
$K$ & Opciones por ranura; $K = 8$ \\
$I$ & Conjunto de instancias de búsqueda; $n = |I| = 100$ \\
$c$ & Configuración; $c[j]$ es la opción elegida en la ranura $j$ \\
$A(c) \subseteq I$ & Instancias que $c$ resuelve (cumplen todas sus restricciones) \\
$F(c) = I \setminus A(c)$ & Instancias que $c$ falla \\
$v_c(i)$ & Familia violada por $c$ en la instancia $i \in F(c)$, según la traza \\
$\Acc(c) = |A(c)|/n$ & Precisión de $c$ sobre la búsqueda \\
$V_j$ & Opciones vivas de la ranura $j$ (no podadas) \\
$R$, $P$ & Número de reinicios y máximo de pasadas por reinicio \\
$L(o,b)$, $G(o,b)$ & Instancias que $o$ pierde y gana frente a $b$ (sección~\ref{sec:poda}) \\
$D(j,o)$ & Vecindarios en los que $o$ fue derrotada en la ranura $j$ \\
$\tau$, $\kappa$, $\mu$ & Pérdida neta mínima para una derrota, vecindarios con derrota necesarios para podar y mínimo de opciones vivas por ranura \\
\bottomrule
\end{tabularx}
\end{table}

\subsection{Vecindad por coordenadas}
Dada una configuración $c$ y una ranura $j$, la vecindad de $c$ en la coordenada $j$ es el conjunto de configuraciones que difieren de $c$ solo en esa ranura y cuya opción está viva o es la opción vigente:
\[
N_j(c) = \{\, c' : c'[h] = c[h] \ \text{para todo}\ h \neq j,\ \ c'[j] \in V_j \cup \{c[j]\} \,\}.
\]
Un paso del método evalúa toda la vecindad $N_j(c)$ y reemplaza $c$ por el elemento de mayor precisión. Una pasada aplica este paso a las ranuras en el orden rol, estrategia, formato, estilo y cierre.

\subsection{Criterio de poda basado en trazas}\label{sec:poda}
Sea $b$ la opción ganadora de la ranura $j$ y $o$ otra opción evaluada en el mismo paso. Las configuraciones $c_b$ y $c_o$ coinciden en todas las ranuras excepto en $j$, de modo que cualquier diferencia entre sus resultados se atribuye al cambio de opción. Se definen las instancias perdidas y ganadas por $o$ frente a $b$:
\[
L(o,b) = A(c_b) \cap F(c_o), \qquad G(o,b) = A(c_o) \cap F(c_b),
\]
y la pérdida neta
\[
\Delta(o,b) = |L(o,b)| - |G(o,b)|.
\]
Se llama vecindario del paso a la tupla $u$ de índices de las otras cuatro ranuras. Si $\Delta(o,b) \geq \tau$, se registra una derrota de $o$ en el vecindario $u$. Sea $D(j,o)$ el conjunto de vecindarios distintos en los que $o$ ha sido derrotada en la ranura $j$. La opción $o$ se elimina de $V_j$ cuando se cumplen las dos condiciones siguientes:
\[
|D(j,o)| \geq \kappa \qquad \text{y} \qquad |V_j| > \mu.
\]
Exigir derrotas en $\kappa$ vecindarios distintos responde a la interacción entre ranuras: una opción que pierde con un vecindario puede ganar con otro, de modo que una sola observación no es evidencia suficiente. Con $\kappa = 2$ la opción debe perder en dos contextos diferentes, que en la práctica aparecen en pasadas o reinicios distintos. El mínimo $\mu = K/2 = 4$ garantiza que la poda nunca descarte más de la mitad de las opciones de una ranura.

Para cada poda se registra, por cada vecindario de $D(j,o)$, la opción ganadora, las precisiones y los conjuntos $L$ y $G$; además, el multiconjunto de familias
\[
\Phi(o) = \{\, v_{c_o}(i) : i \in L(o,b) \,\}
\]
acumulado sobre esos vecindarios, que indica qué tipo de restricción rompe la opción, y un extracto de la salida del modelo en una instancia perdida. Esa información justifica la decisión y se guarda en poda.json; las derrotas que todavía no produjeron una poda se guardan en derrotas.csv. Las opciones candidatas se procesan en orden decreciente de pérdida neta acumulada, de modo que, si el mínimo $\mu$ impide podar todas, se eliminan primero las que más daño causan.

La comparación se hace instancia por instancia y no solo por precisión agregada. Dos opciones con la misma precisión pueden fallar instancias distintas; en ese caso la pérdida neta es cero y el criterio no las poda sin evidencia.

\subsection{Pseudocódigo}

\begin{algorithm}[H]
\caption{Ascenso por coordenadas con poda guiada por trazas.}
\begin{algorithmic}[1]
\Require puntos iniciales $c_1, \dots, c_R$; parámetros $P$, $\tau$, $\kappa$, $\mu$
\Ensure mejor configuración observada y su precisión
\State $V_j \gets \{0, \dots, K-1\}$ para cada ranura $j$
\For{$r = 1$ \textbf{hasta} $R$}
  \State $c \gets c_r$; \ $p \gets \Acc(c)$
  \For{$\text{pasada} = 1$ \textbf{hasta} $P$}
    \State $\text{mejoró} \gets$ falso
    \ForAll{ranura $j$ en (rol, estrategia, formato, estilo, cierre)}
      \State evaluar $\Acc(c')$ y las trazas de cada $c' \in N_j(c)$
      \State $b \gets$ opción de mayor precisión; si empata con $c[j]$, $b \gets c[j]$
      \State \Call{Podar}{$j, b, c$}
      \If{$\Acc(c \text{ con } j = b) > p$} $\text{mejoró} \gets$ verdadero \EndIf
      \State $c[j] \gets b$; \ $p \gets \Acc(c)$
    \EndFor
    \If{\textbf{no} mejoró} terminar el reinicio $r$ \EndIf
  \EndFor
  \State guardar $(c, p)$ como resultado del reinicio $r$
\EndFor
\State \Return el resultado de mayor precisión entre los $R$ reinicios
\end{algorithmic}
\end{algorithm}

\begin{algorithm}[H]
\caption{Procedimiento \textsc{Podar}$(j, b, c)$.}
\begin{algorithmic}[1]
\State $u \gets$ índices de $c$ en las ranuras distintas de $j$
\ForAll{opción $o \neq b$ evaluada en el paso}
  \If{$|L(o,b)| - |G(o,b)| \geq \tau$} $D(j,o) \gets D(j,o) \cup \{u\}$ \EndIf
\EndFor
\State $Q \gets \{\, o \in V_j : o \neq b,\ |D(j,o)| \geq \kappa \,\}$, ordenado de mayor a menor pérdida neta acumulada
\ForAll{$o \in Q$}
  \If{$|V_j| \leq \mu$} terminar \EndIf
  \State $V_j \gets V_j \setminus \{o\}$
  \State registrar $o$, la evidencia de cada vecindario de $D(j,o)$, $\Phi(o)$ y un extracto de la salida
\EndFor
\end{algorithmic}
\end{algorithm}

Todas las evaluaciones pasan por una memoria (memo) indexada por la configuración. Si una configuración ya se evaluó en un paso, una pasada o un reinicio anteriores, se reutiliza su resultado sin una nueva consulta al oráculo.

\subsection{Desempate, puntos iniciales y parada}
Si varias opciones alcanzan la precisión máxima y la opción vigente está entre ellas, se conserva la vigente; en otro caso se elige la de menor índice. Así la configuración solo cambia cuando hay evidencia de mejora y el resultado no depende del orden de evaluación.

Los puntos iniciales son la configuración con todos los índices en 0 y $R-1$ configuraciones generadas con random.Random(SEMILLA). Ninguno usa información previa sobre qué opción es preferible.

Un reinicio termina cuando una pasada completa no mejora la precisión o cuando se alcanzan $P$ pasadas. La búsqueda termina al completar los $R$ reinicios.

\subsection{Propiedades}\label{sec:propiedades}

\begin{proposicion}[No decrecimiento]\label{prop:monotonia}
Dentro de un reinicio, la precisión de la configuración vigente nunca disminuye.
\end{proposicion}
\begin{proof}
En cada paso, la configuración vigente pertenece a $N_j(c)$, porque $c[j]$ se incluye siempre en la vecindad, aunque haya sido podada. La opción elegida maximiza la precisión sobre ese conjunto, luego $\Acc(c \text{ con } j = b) \geq \Acc(c)$. Con temperatura 0.0 el oráculo es determinista y la precisión de $c$ no cambia entre consultas.
\end{proof}

\begin{proposicion}[Terminación y cota de consultas]\label{prop:cota}
Cada reinicio realiza a lo sumo $1 + P \cdot m \cdot K$ evaluaciones y la búsqueda completa a lo sumo $R\,(1 + P \cdot m \cdot K)$.
\end{proposicion}
\begin{proof}
Cada reinicio hace una evaluación inicial y como máximo $P$ pasadas. En cada pasada, cada una de las $m$ ranuras evalúa $|V_j \cup \{c[j]\}| \leq K$ configuraciones. Sumando, el reinicio realiza a lo sumo $1 + P \cdot m \cdot K$ evaluaciones, y la memoria solo puede reducir este número. Como $R$, $P$, $m$ y $K$ son finitos, el algoritmo termina.
\end{proof}

\begin{proposicion}[Mínimo de opciones]\label{prop:minimo}
En todo momento $|V_j| \geq \mu$ para cada ranura $j$.
\end{proposicion}
\begin{proof}
Inicialmente $|V_j| = K \geq \mu$. \textsc{Podar} solo elimina una opción cuando $|V_j| > \mu$, y elimina una por vez, por lo que después de cada eliminación $|V_j| \geq \mu$.
\end{proof}

\begin{proposicion}[Reproducibilidad]\label{prop:reproducible}
Con los mismos datos, modelo, semilla y parámetros, dos ejecuciones producen la misma secuencia de configuraciones evaluadas y la misma configuración final.
\end{proposicion}
\begin{proof}
Con temperatura 0.0 la decodificación es voraz (greedy), así que la respuesta a cada par de configuración e instancia es fija. Los puntos iniciales dependen solo de SEMILLA, el orden de las ranuras es fijo y el desempate es determinista. Por inducción sobre los pasos, todas las decisiones coinciden.
\end{proof}

Las proposiciones no garantizan optimalidad global. El ascenso por coordenadas puede detenerse en una configuración en la que ningún cambio de una sola ranura mejora, aunque exista una mejor que exija cambiar dos ranuras a la vez. La poda agrega un segundo riesgo: una opción descartada en un contexto podría ser la mejor en otro, porque las ranuras interactúan. Los parámetros $\tau$, $\kappa$ y $\mu$, y la reevaluación de la opción vigente aunque esté podada, limitan ese riesgo sin eliminarlo.

\subsection{Análisis de complejidad}
Se distinguen tres costos: el número de consultas al oráculo, el costo de cada consulta y el trabajo propio de la heurística.

\subsubsection{Número de consultas}
Sea $Q$ el número de configuraciones distintas evaluadas. Por la proposición~\ref{prop:cota}:
\[
Q \leq R\,(1 + P \cdot m \cdot K) = 3 \cdot (1 + 3 \cdot 5 \cdot 8) = 363.
\]
Una vez que hay podas, cada pasada evalúa $\sum_j |V_j \cup \{c[j]\}|$ configuraciones, con $\mu \leq |V_j| \leq K$. Si todas las ranuras llegan a $\mu = 4$ opciones vivas, una pasada cuesta como máximo $m(\mu + 1) = 25$ evaluaciones en lugar de $m \cdot K = 40$.

\subsubsection{Costo de una consulta}
Una consulta genera una respuesta por instancia no guardada en caché y verifica sus restricciones. Si $t_g$ es el tiempo de generar una respuesta de a lo sumo 384 tokens nuevos y $t_v$ el de verificar una restricción, y cada instancia tiene a lo sumo 5 restricciones, el costo de una consulta es
\[
O\big(n\,(t_g + 5\,t_v)\big) = O(n \cdot t_g),
\]
porque la generación domina. El oráculo agrupa las peticiones en lotes de hasta 16 prompts, lo que reduce la constante pero no el orden.

\subsubsection{Trabajo de la heurística}
\textsc{Podar} calcula $L(o,b)$ y $G(o,b)$ para a lo sumo $K-1$ opciones con conjuntos de identificadores, en $O(n)$ cada una, y ordena a lo sumo $K$ candidatas. La pérdida neta acumulada se suma sobre a lo sumo $R \cdot P$ vecindarios por opción. El costo por paso es $O(K \cdot n + K \cdot R \cdot P + K \log K) = O(K \cdot n)$, porque $R \cdot P \leq n$ con los valores del proyecto, y en toda la búsqueda $O(R \cdot P \cdot m \cdot K \cdot n)$, despreciable frente a la generación.

\subsubsection{Costo total y comparación}
El tiempo total es
\[
T = O(Q \cdot n \cdot t_g) = O(R \cdot P \cdot m \cdot K \cdot n \cdot t_g),
\]
polinomial en los parámetros. La enumeración exhaustiva evalúa las $K^m$ configuraciones, con un costo $O(K^m \cdot n \cdot t_g)$, exponencial en $m$. La tabla~\ref{tab:comparacion} compara ambas cantidades con los valores del proyecto.

\begin{table}[H]
\caption{Generaciones del modelo en el peor caso.}\label{tab:comparacion}
\begin{tabularx}{\textwidth}{@{}Lp{4.3cm}p{3.8cm}@{}}
\toprule
Método & Configuraciones & Generaciones ($n = 100$) \\
\midrule
Enumeración exhaustiva (temperatura 0.0) & $8^5 = 32\,768$ & $3\,276\,800$ \\
Ascenso por coordenadas, $R = 3$, $P = 3$ & $\leq 363$ & $\leq 36\,300$ \\
Ascenso con poda, pasadas 2 y 3 con $\mu = 4$ & $\leq 3\,(1 + 40 + 2 \cdot 25) = 273$ & $\leq 27\,300$ \\
\bottomrule
\end{tabularx}
\end{table}

La última fila es una cota del caso en el que todas las ranuras alcanzan $\mu$ opciones vivas al final de la primera pasada de cada reinicio; el número real depende de cuántas podas ocurran y de cuántas configuraciones se repitan entre reinicios.

\subsubsection{Memoria}
La memoria guarda, por configuración evaluada, la precisión y las trazas de las instancias falladas, cada una con una salida de a lo sumo 384 tokens. El espacio es $O(Q \cdot n \cdot \ell)$, con $\ell$ la longitud máxima de una salida. La caché del oráculo en disco crece en la misma proporción y por eso se exporta comprimida.

\subsection{Estructuras de datos}\label{sec:estructuras}

\begin{table}[H]
\caption{Estructuras de datos del notebook de búsqueda.}
\begin{tabularx}{\textwidth}{@{}p{2.2cm}p{3cm}L@{}}
\toprule
Estructura & Tipo & Contenido y uso \\
\midrule
memo & diccionario & Llave: tupla de los cinco índices y la temperatura. Valor: precisión, aciertos, diccionario de id a familia violada, contador de familias y trazas completas. Evita repetir consultas y permite comparar opciones instancia por instancia. \\
registro & lista de diccionarios & Una entrada por consulta nueva, con su contexto (reinicio, pasada, ranura, opción) y su resultado. Se exporta a evaluaciones.csv y registro\_busqueda.json. \\
vivas & diccionario de listas & Conjunto $V_j$ de opciones vivas de cada ranura. \\
derrotas & diccionario & Llave: (ranura, opción). Valor: diccionario de vecindario a evidencia (opción ganadora, precisiones, instancias perdidas y ganadas, familias). Implementa $D(j,o)$. Se exporta a derrotas.csv. \\
podas & lista de diccionarios & Decisiones de poda con su evidencia. Se exporta a poda.json. \\
historial & lista de reales & Precisión de cada consulta nueva, en orden, para la curva de búsqueda. \\
\bottomrule
\end{tabularx}
\end{table}

% ---------------------------------------------------------------- 6
\section{Procedimiento de ejecución}

\subsection{Notebook de búsqueda}
\begin{enumerate}[leftmargin=*]
\item Abrir Google Colab, elegir Archivo, Abrir notebook, GitHub, y pegar la dirección del repositorio. Seleccionar notebooks/01\_busqueda.ipynb.
\item En Entorno de ejecución, Cambiar tipo de entorno de ejecución, seleccionar GPU T4.
\item Ejecutar la celda de la sección 1. Al terminar, elegir Entorno de ejecución, Reiniciar sesión.
\item Ejecutar la sección 2. Si USAR\_DRIVE es verdadero, autorizar el acceso a Google Drive; la caché se guarda en MyDrive/ADA\_Oraculo\_G01/cache\_oraculo.json. Si existe una caché de una ejecución anterior con el mismo modelo y semilla, copiarla en esa ruta antes de ejecutar: las configuraciones ya generadas se resuelven sin usar la GPU.
\item Ejecutar las secciones 3 y 4. La sección 4 imprime una línea por consulta y una línea por cada poda. Con la GPU gratuita de Colab la búsqueda completa puede tardar varias horas.
\item Ejecutar las secciones 5 a 7: tablas de evaluaciones, análisis de trazas y validación con oraculo.validar sobre las 87 instancias reservadas.
\item Ejecutar la sección 8. Se escriben los archivos de resultados/, entrega.json y se descarga resultados\_busqueda.zip. Descargar también el notebook con sus salidas (Archivo, Descargar, .ipynb) y subir todo al repositorio.
\end{enumerate}

\subsection{Notebook de evaluación}
\begin{enumerate}[leftmargin=*]
\item Si se cuenta con un archivo de prueba propio, subirlo al panel de archivos de Colab con el nombre datos\_prueba.json, en /content. También se acepta una ruta en la variable de entorno ORACULO\_DATOS o el archivo MyDrive/datos\_prueba.json.
\item Abrir notebooks/02\_evaluacion.ipynb, elegir GPU T4, ejecutar la sección 1 y reiniciar la sesión.
\item Ejecutar las secciones 2 a 4 sin modificar el código. La salida indica el origen de los datos, la precisión, el número de aciertos y las familias de restricción que más fallaron. Si no se encuentra datos\_prueba.json, se evalúa sobre la partición de validación y la salida lo indica.
\end{enumerate}

\subsection{Reanudación después de una desconexión}\label{sec:reanudacion}
Si Colab se desconecta durante la búsqueda, basta con abrir de nuevo el notebook, ejecutar la sección 1, reiniciar y ejecutar desde la sección 2. El oráculo carga la caché de Google Drive y las configuraciones ya generadas se resuelven sin usar la GPU. Por la proposición~\ref{prop:reproducible}, la búsqueda recorre las mismas configuraciones en el mismo orden y continúa desde el punto donde se interrumpió.

\subsection{Parámetros}

\begin{table}[H]
\caption{Parámetros de la búsqueda (sección 2 del notebook).}
\begin{tabularx}{\textwidth}{@{}p{4.2cm}p{2.2cm}L@{}}
\toprule
Parámetro & Valor & Efecto \\
\midrule
GRUPO & "G01" & Identificador escrito en entrega.json \\
SEMANA & 5 & Semana académica escrita en entrega.json; debe confirmarse con la docente \\
MODELO & "qwen17b" & Modelo usado por el oráculo; cambiarlo invalida la caché \\
TEMPERATURA & 0.0 & Temperatura de todas las configuraciones evaluadas \\
SEMILLA & 1 & Semilla del oráculo y de los puntos iniciales aleatorios \\
MAX\_PASADAS ($P$) & 3 & Pasadas máximas por reinicio \\
N\_REINICIOS ($R$) & 3 & Número de puntos iniciales \\
UMBRAL\_PODA ($\tau$) & 3 & Pérdida neta mínima, en instancias, para registrar una derrota \\
CONTEXTOS\_PODA ($\kappa$) & 2 & Vecindarios distintos con derrota necesarios para podar una opción \\
MIN\_VIVAS ($\mu$) & 4 & Opciones mínimas que conserva cada ranura \\
USAR\_DRIVE & True & Guarda la caché en Google Drive \\
\bottomrule
\end{tabularx}
\end{table}

% ---------------------------------------------------------------- 7
\section{Interpretación de las salidas}

\begin{table}[H]
\caption{Salidas del sistema y su significado.}
\begin{tabularx}{\textwidth}{@{}p{5.6cm}L@{}}
\toprule
Salida & Significado \\
\midrule
consulta k \ p \ mejor q \ ranura=o \ s & Consulta nueva número k: precisión p de la configuración evaluada, mejor precisión q hasta ese momento, opción probada y segundos empleados \\
poda ranura=o: derrotada en k contextos, pérdida neta x\dots & Se descartó la opción o porque perdió en k vecindarios distintos, con x instancias netas en total; se muestran las familias más frecuentes en las instancias perdidas \\
fin pasada p (reinicio r) & Precisión y configuración vigente al terminar la pasada \\
Resultado(precision=\dots, fallos=f/n) & Representación de un resultado del oráculo \\
validación: p (87 inst) & Precisión en la partición de validación seguida de las familias que más fallaron \\
consultas que habría hecho el ascenso sin poda & Cota superior para el mismo número de pasadas: cuenta la opción vigente en cada paso y no descuenta las configuraciones repetidas. El ahorro real se mide con configuraciones distintas (sección~\ref{sec:conpoda}) \\
entrega.json & Configuración final: grupo, semana, índices de las cinco ranuras y temperatura \\
\bottomrule
\end{tabularx}
\end{table}

La precisión de búsqueda y la de validación miden cosas distintas. La primera se obtuvo sobre las mismas instancias usadas para elegir la configuración y tiende a sobrestimar el desempeño; la segunda se mide sobre instancias no usadas en la búsqueda. Ninguna de las dos sustituye la evaluación oficial de la docente sobre datos privados.

% ---------------------------------------------------------------- 8
\section{Verificación}

\subsection{Comprobaciones internas}
\begin{itemize}[leftmargin=*]
\item El notebook de búsqueda verifica que la partición de búsqueda tenga 100 instancias, la de validación 87 y que cada ranura tenga 8 opciones.
\item El notebook de evaluación verifica que entrega.json contenga exactamente las cinco ranuras y la temperatura, que cada índice sea un entero entre 0 y 7 y que la temperatura sea un valor permitido. Si alguna condición falla, se detiene con un mensaje que identifica el campo.
\end{itemize}

\subsection{Prueba con oráculo sintético}
El archivo pruebas/prueba\_busqueda\_simulada.py ejecuta las funciones del notebook de búsqueda sin GPU ni dependencias externas. Reemplaza el oráculo por uno sintético y determinista con una configuración óptima conocida: cada ranura distinta de la óptima hace fallar un subconjunto fijo de instancias. Se ejecuta desde la raíz del repositorio con el comando python pruebas/prueba\_busqueda\_simulada.py.

\begin{table}[H]
\caption{Comprobaciones de la prueba sintética.}
\begin{tabularx}{\textwidth}{@{}Lp{4.6cm}p{2cm}@{}}
\toprule
Comprobación & Propiedad verificada & Resultado \\
\midrule
Se encuentra la configuración óptima y la precisión reportada es 100\,\% & Corrección en un caso con óptimo conocido & Correcta \\
Ninguna configuración se consulta dos veces & Funcionamiento de la memoria & Correcta \\
Cada consulta nueva queda en el registro & Trazabilidad & Correcta \\
Ningún reinicio termina por debajo de su punto inicial & Proposición~\ref{prop:monotonia} & Correcta \\
Toda ranura conserva al menos $\mu$ opciones & Proposición~\ref{prop:minimo} & Correcta \\
Toda poda se apoya en derrotas en al menos $\kappa$ vecindarios & Criterio de la sección~\ref{sec:poda} & Correcta \\
Cada derrota tiene una pérdida neta de al menos $\tau$ instancias & Criterio de la sección~\ref{sec:poda} & Correcta \\
Las consultas no superan 363 & Proposición~\ref{prop:cota} & Correcta \\
\bottomrule
\end{tabularx}
\end{table}

En la ejecución de la prueba, la búsqueda realizó 89 consultas y 16 podas, frente a la cota de 363 consultas sin poda. El programa termina con código 0 si todas las comprobaciones se cumplen y con código 1 en caso contrario.

\subsection{Comparación con y sin poda}\label{sec:conpoda}
El archivo pruebas/comparar\_poda\_sintetica.py ejecuta el mismo código sobre 50 configuraciones óptimas elegidas con la semilla 2026, una vez con la poda activada y otra con un umbral de poda inalcanzable, que equivale a no podar. Se ejecuta con python pruebas/comparar\_poda\_sintetica.py y debe imprimir los valores de la tabla~\ref{tab:conpoda}.

\begin{table}[H]
\caption{Consultas al oráculo sintético con y sin poda (50 configuraciones óptimas).}\label{tab:conpoda}
\begin{tabularx}{\textwidth}{@{}Lp{4.2cm}p{3cm}@{}}
\toprule
Variante & Consultas: media (mínimo--máximo) & Encuentra el óptimo \\
\midrule
Sin poda, 2 reinicios & 90.0 (78--92) & 50 de 50 \\
Con poda, 2 reinicios & 76.4 (72--78) & 50 de 50 \\
Sin poda, 3 reinicios & 116.6 (106--120) & 50 de 50 \\
Con poda, 3 reinicios & 89.4 (84--92) & 50 de 50 \\
\bottomrule
\end{tabularx}
\end{table}

Con los mismos reinicios, la poda reduce las consultas entre un 15\,\% y un 23\,\% sin perder el óptimo. En este oráculo cada ranura afecta a las instancias de forma independiente, así que la prueba comprueba la corrección y el ahorro de la poda, no la calidad de la solución en el problema real. El informe compara además las dos versiones sobre los datos del curso y sobre los datos propios. El archivo pruebas/grafica\_ahorro\_poda.py resume los tres escenarios en resultados/ahorro\_poda.png.

\subsection{Datos de prueba propios}\label{sec:datospropios}
La carpeta datos de prueba contiene datos\_visibles\_100.json, un conjunto de 100 instancias construido por el grupo con el formato de la sección~\ref{sec:instancia}: 40 de búsqueda, 25 de validación y 35 descartadas, con 45 preguntas propias y restricciones de 26 familias de IFEvalG. Sus instancias tienen menos restricciones que las del curso (1.68 en promedio en búsqueda, frente a 2.45), por lo que sus precisiones son más altas y no se comparan con las de datos\_visibles.json.

Este archivo también sirve para probar el camino que seguirá la docente con sus datos privados: se copia con el nombre datos\_prueba.json en el panel de archivos de Colab (/content) y se ejecuta 02\_evaluacion.ipynb desde la sección 2. El notebook debe informar /content/datos\_prueba.json como origen y evaluar 65 instancias, porque omite las marcadas como descartadas.

\subsection{Prueba mínima con el modelo}
Para comprobar la instalación sin ejecutar la búsqueda completa, después de la sección 2 del notebook de búsqueda se puede evaluar una configuración sobre cinco instancias:
\begin{center}
r = oraculo.evaluar(config, busqueda[:5], semilla=1)
\end{center}
El resultado es correcto si r.n es 5, r.precision está en $[0,1]$, el número de trazas es $5 - 5 \cdot \text{r.precision}$ y cada traza tiene los campos id, violo y salida.

\subsection{Reproducibilidad de los resultados}
La configuración de entrega.json debe coincidir con la mejor configuración del notebook de búsqueda ejecutado y con el campo mejor\_config de registro\_busqueda.json. La precisión de búsqueda puede recalcularse a partir de evaluaciones.csv, y la caché comprimida permite repetir el análisis de trazas sin GPU.

La ejecución entregada debe reproducir estos valores: configuración rol~=~1, estrategia~=~2, formato~=~3, estilo~=~5, cierre~=~1 y temperatura~0.0; precisión de búsqueda 18.00\,\% (18 de 100); precisión de validación 12.64\,\% (11 de 87); 140 configuraciones evaluadas y 20 opciones podadas. El notebook de evaluación, sin datos\_prueba.json, debe obtener la misma precisión de validación.

% ---------------------------------------------------------------- 9
\section{Supuestos, limitaciones y trabajo pendiente}

\subsection{Supuestos}
\begin{itemize}[leftmargin=*]
\item La temperatura es 0.0 y la decodificación es voraz, de modo que el oráculo se comporta como una función determinista de la configuración y la instancia.
\item La búsqueda usa solo las 100 instancias de búsqueda. La partición de validación no interviene en la elección de la configuración.
\item Los datos privados de prueba tienen el formato de la sección~\ref{sec:instancia} y sus familias de restricción están en el registro IFEvalG de open-instruct.
\end{itemize}

\subsection{Limitaciones y errores conocidos}
\begin{itemize}[leftmargin=*]
\item La búsqueda es heurística y puede terminar en un óptimo local (sección~\ref{sec:propiedades}).
\item La poda puede descartar una opción que sería la mejor en otro contexto de las demás ranuras.
\item Con $n = 100$ cada instancia representa un punto porcentual. Diferencias de una o dos instancias entre configuraciones pueden deberse a la muestra y no a una ventaja real.
\item Una ejecución completa toma varias horas en una T4 gratuita y puede verse interrumpida por desconexiones de Colab.
\item Si los datos de prueba incluyen familias que no están en IFEvalG, el oráculo lanza KeyError salvo que se le pasen registros adicionales mediante el parámetro registros de Oraculo.
\item Si una instancia no trae el campo restriccion, la opción de cierre que repite las restricciones no agrega texto al prompt.
\item El notebook de evaluación busca también /content/drive/MyDrive/datos\_prueba.json, pero no monta Google Drive; esa ruta solo se usa si Drive se montó antes en la sesión. Las rutas /content/datos\_prueba.json y ORACULO\_DATOS no tienen esa restricción.
\end{itemize}

\subsection{Componentes no desarrollados}\label{sec:pendiente}
\begin{itemize}[leftmargin=*]
\item Estrategias de reparto del presupuesto de consultas (uniforme y adaptativo), Successive Halving y asignación por UCB1.
\item Exploración de las temperaturas 0.3 y 0.7.
\item Medición de tiempos con tamaños crecientes de $n$ y comparación con la función de eficiencia teórica.
\end{itemize}

% ---------------------------------------------------------------- referencias
\section*{Referencias}
\addcontentsline{toc}{section}{Referencias}
\begingroup
\setlength{\parindent}{-1.27cm}
\setlength{\leftskip}{1.27cm}
\raggedright

Allen Institute for AI. (2025). \textit{open-instruct} [Repositorio de software]. GitHub. \url{https://github.com/allenai/open-instruct}

Melo, D. (2026). \textit{Proyecto-Oraculo-AyD} [Repositorio de software]. GitHub. \url{https://github.com/DanielMelo404/Proyecto-Oraculo-AyD}

Qwen Team. (2025). \textit{Qwen3 technical report}. arXiv. \url{https://arxiv.org/abs/2505.09388}

Wright, S. J. (2015). Coordinate descent algorithms. \textit{Mathematical Programming, 151}(1), 3--34. \url{https://doi.org/10.1007/s10107-015-0892-3}

Zhou, J., Lu, T., Mishra, S., Brahma, S., Basu, S., Luan, Y., Zhou, D., y Hou, L. (2023). \textit{Instruction-following evaluation for large language models}. arXiv. \url{https://arxiv.org/abs/2311.07911}

\endgroup

\end{document}
